\BeleskeID{01-s031}
% element: 01-s031:e01
% element: 01-s031:e02
% element: 01-s031:e03
% element: 01-s031:e04

% element: 01-s031:d01
NI i NILI mogu se čitati kao uslovi I i ILI čije je ispunjenje označeno nulom. Negacija menja vrednost kojom se uslov iskazuje, dok kombinacija ulaznih vrednosti koja ispunjava taj uslov ostaje ista.

Pri kaskadiranju je presudno šta sledeći stepen prima na svom ulazu. Ako prvi NI stepen daje aktivnu nulu za ispunjen uslov, drugi običan NI ulaz tu nulu ne tumači kao jedinicu. Zato dva NI kola ne mogu neposredno zameniti kaskadu I kola uz samo završnu negaciju. Za svaki međusignal treba napisati njegovu stvarnu logičku vrednost, umesto zaključivati samo iz naziva kola.
